\documentclass[11pt]{article}
\usepackage{amsmath,amssymb,amsthm,geometry}
\geometry{margin=1in}
\title{Synchronous Local Path Broadcast and the Limits of AC$^0$ Normalization}
\author{Inacio F. Vasquez}
\date{Corrected scope: November 26, 2026}
\newtheorem{theorem}{Theorem}
\begin{document}
\maketitle
\begin{abstract}
We distinguish bounded per-output circuit dependence from bounded aggregate
observation capacity. The former does not imply the latter. We prove an
exact minimum-round bound for seeded broadcast on a path under synchronous
nearest-neighbor communication and distributed output requirements.
The historical unrestricted AC$^0$ normalization and universal
EntropyDepth separation are not established.
\end{abstract}
\section{Parallel update counterexample}
Let $X$ be uniform on $\{0,1\}^n$, and let $Y_i=1-X_i$.
Each of the $n$ outputs is computed by a depth-one fan-in-one gate.
Relative to an input-independent initial observation, $Y$ exposes
$I(X;Y)=H(X)=n$ bits. Thus bounded fan-in, arity and depth do not
supply a constant aggregate information bound. If the initial
observation already contains $X$, the conditional information gain is
zero instead; the observation model must be fixed.
A bounded dependency cone per output also need not be a geometrically
local neighborhood of the original constraint graph.
\section{Exact restricted distributed result}
\begin{theorem}
Let $P_n$ be the path with $n\geq1$ vertices.
Only the first vertex receives an arbitrary bit $b$ initially, with
all other initial data fixed independently of $b$.
Each deterministic synchronous round updates each vertex using only
its own and adjacent states from the preceding round.
Require every vertex to output $b$ correctly for both seeds, with
no nonlocal communication, input-dependent auxiliary initialization,
or substitute global readout.
The minimum worst-case number of rounds is $n-1$.
\end{theorem}
\begin{proof}
Induct on $t$ to show the state of each vertex after $t$ rounds
depends only on the initial states at distance at most $t$.
At time zero this holds by definition; the union of the preceding
neighbor dependency balls is contained in the next ball.
If $t<n-1$, the last vertex has identical states for the two seeds,
but different required outputs. Hence no correct algorithm can
finish by that round for both seeds.
Broadcasting the known bit to adjacent vertices each round attains
$n-1$ rounds, including zero rounds for $n=1$.
\end{proof}
\section{Scope and open obligations}
This theorem concerns distributed output and round complexity,
not deciding satisfiability of an already known family.
The unseeded implication chain has both all-zero and all-one
satisfying assignments; unit propagation does not start without a seed.
With the unit clause $x_1=1$, synchronous unit propagation forces
the first $\min(n,t+1)$ variables after $t$ rounds.
This is an algorithm-specific result.
No universal $\mathrm{EDA}_A(F)=\Theta(n)$ assertion follows.
The proposed B** simulation requires explicit encoding, access,
aggregate capacity, terminal uncertainty and cost preservation.
Bounded quantifier rank alone is not strict neighborhood locality.
TRAP remains an unproved assumption; no general SAT or P versus NP
lower bound is derived. Lean verification of the path theorem is
not claimed.
\end{document}
